\documentclass[12pt,oneside]{book}

% =========================================================
% PAQUETES
% =========================================================
\usepackage[utf8]{inputenc}
\usepackage[T1]{fontenc}
\usepackage[spanish,es-nodecimaldot]{babel}

\usepackage[
    top=2.5cm,
    bottom=2.5cm,
    left=3cm,
    right=2.5cm,
    headheight=15pt
]{geometry}

\usepackage{amsmath}
\usepackage{amssymb}
\usepackage{mathtools}

\usepackage{graphicx}
\usepackage{float}
\usepackage{caption}
\usepackage{subcaption}

\usepackage{booktabs}
\usepackage{array}
\usepackage{longtable}
\usepackage{tabularx}

\usepackage{enumitem}

\usepackage{xcolor}
\usepackage[most]{tcolorbox}

\usepackage{fancyhdr}
\usepackage{ragged2e}

% =========================================================
% RUTAS DE IMÁGENES
% =========================================================
\graphicspath{
    {./img/}
    {./graficos/}
    {./figuras/}
}

% =========================================================
% METADATOS DEL DOCUMENTO
% =========================================================
\newcommand{\universidad}{Universidad de Buenos Aires (UBA)}
\newcommand{\carrera}{Carrera de Abogacía --- Ciudad Autónoma de Buenos Aires, Argentina}
\newcommand{\materia}{Lectocomprensión de Textos Jurídicos en Inglés}
\newcommand{\materiacorta}{Inglés Jurídico}
\newcommand{\titulo}{Apuntes de clase: estudio del texto \emph{``What is the UK Constitution?''}}
\newcommand{\autor}{Isaac Edgar Camacho Ocampo}
\newcommand{\anio}{2026}

\usepackage[
    colorlinks=true,
    linkcolor=blue!50!black,
    citecolor=green!40!black,
    urlcolor=blue!60!black,
    pdftitle={Apuntes de clase: estudio del texto "What is the UK Constitution?"},
    pdfauthor={Isaac Edgar Camacho Ocampo},
    pdfsubject={Apuntes universitarios},
    bookmarks=true
]{hyperref}

% =========================================================
% ÍNDICE
% =========================================================
\setcounter{tocdepth}{2}

% =========================================================
% CONFIGURACIÓN DE PÁRRAFOS Y LISTAS
% =========================================================
\setlength{\parindent}{0pt}
\setlength{\parskip}{0.5em}

\setlist[itemize]{
    topsep=0.3em,
    itemsep=0.2em
}

\setlist[enumerate]{
    topsep=0.3em,
    itemsep=0.2em
}

\clubpenalty=10000
\widowpenalty=10000

% =========================================================
% ENCABEZADOS Y PIES DE PÁGINA
% =========================================================
\pagestyle{fancy}
\fancyhf{}

\fancyhead[L]{\nouppercase{\leftmark}}
\fancyhead[R]{\materiacorta}

\fancyfoot[C]{\thepage}

\renewcommand{\headrulewidth}{0.4pt}
\renewcommand{\footrulewidth}{0pt}

\fancypagestyle{plain}{
    \fancyhf{}
    \fancyfoot[C]{\thepage}
    \renewcommand{\headrulewidth}{0pt}
}

% =========================================================
% COMANDOS MATEMÁTICOS
% =========================================================
\newcommand{\abs}[1]{\left|#1\right|}
\newcommand{\norm}[1]{\left\lVert#1\right\rVert}
\newcommand{\vect}[1]{\mathbf{#1}}
\newcommand{\deriv}[2]{\frac{d#1}{d#2}}
\newcommand{\pderiv}[2]{\frac{\partial#1}{\partial#2}}

% =========================================================
% CAJAS ACADÉMICAS
% =========================================================
\newtcolorbox{definition}[1]{
    enhanced,
    breakable,
    title={Definición: #1},
    fonttitle=\bfseries,
    colback=blue!3,
    colframe=blue!50!black,
    boxrule=0.8pt,
    arc=2mm
}

\newtcolorbox{important}{
    enhanced,
    breakable,
    title={Importante},
    fonttitle=\bfseries,
    colback=yellow!8,
    colframe=orange!70!black,
    boxrule=0.8pt,
    arc=2mm
}

\newtcolorbox{example}[1]{
    enhanced,
    breakable,
    title={Ejemplo: #1},
    fonttitle=\bfseries,
    colback=green!4,
    colframe=green!40!black,
    boxrule=0.8pt,
    arc=2mm
}

\newtcolorbox{formula}[1]{
    enhanced,
    breakable,
    title={Fórmula / Regla: #1},
    fonttitle=\bfseries,
    colback=purple!4,
    colframe=purple!50!black,
    boxrule=0.8pt,
    arc=2mm
}

\newtcolorbox{exercise}[1]{
    enhanced,
    breakable,
    title={Ejercicio: #1},
    fonttitle=\bfseries,
    colback=gray!5,
    colframe=gray!60!black,
    boxrule=0.8pt,
    arc=2mm
}

\newtcolorbox{note}{
    enhanced,
    breakable,
    title={Nota},
    fonttitle=\bfseries,
    colback=white,
    colframe=gray!50,
    boxrule=0.6pt,
    arc=2mm
}

% =========================================================
% DOCUMENTO
% =========================================================
\begin{document}

% ---------------------------------------------------------
% PORTADA
% ---------------------------------------------------------
\begin{titlepage}
\thispagestyle{empty}

\begin{center}

\vspace*{1cm}

{\large \textsc{\universidad}}

\vspace{0.2cm}

{\large \carrera}

\vspace{3cm}

{\Huge \textbf{\materia}}

\vspace{1cm}

{\LARGE \titulo}

\vspace{0.8cm}

\textit{Comprender un sistema jurídico distinto del propio es también aprender a mirar el propio con otros ojos.}

\vspace{1.2cm}

\rule{0.70\textwidth}{0.5pt}

\vspace{0.5cm}

{\large Apuntes de estudio}

\vspace{0.5cm}

\rule{0.70\textwidth}{0.5pt}

\vfill

{\large \autor}

\vspace{0.3cm}

{\large \anio}

\end{center}

\vfill

\begin{flushleft}
\textit{Este es un modesto aporte a los estudiantes de ``Lectocomprensión de Textos Jurídicos en Inglés'' no pretende sustituir el material oficial de la materia.}
\end{flushleft}

\end{titlepage}

% ---------------------------------------------------------
% ÍNDICE
% ---------------------------------------------------------
\tableofcontents

% ===========================================================
\chapter{Presentación del texto de trabajo y objetivos de la clase}
% ===========================================================

La clase tuvo como eje el análisis del texto \emph{``What is the UK constitution?''}, empleado como caso de estudio para desarrollar estrategias de lectocomprensión de textos jurídicos en inglés y, al mismo tiempo, para comprender el sistema constitucional del Reino Unido en contraste con el sistema continental codificado, tomando a la Argentina como referencia.

\section{Temas tratados en la clase}

\begin{itemize}
    \item Estrategias de lectocomprensión: identificación de la macroestructura del texto y de los términos transparentes en inglés jurídico.
    \item Estructura y fuentes de la Constitución del Reino Unido: \emph{statutes}, \emph{conventions}, \emph{common law}, derecho de la Unión Europea y derecho internacional.
    \item Falsos cognados jurídicos clave: \emph{jurisprudence}, \emph{statute} y \emph{convention}.
    \item Soberanía parlamentaria, la Corona y los principios centrales del sistema constitucional británico.
    \item Reformas constitucionales históricas y los acuerdos de \emph{devolution} con Escocia, Gales e Irlanda del Norte.
\end{itemize}

\section{Relevancia para la práctica profesional}

\begin{note}
El ejercicio profesional del derecho está cada vez más atravesado por fuentes, contratos y jurisprudencia redactados en inglés: compraventa internacional de mercaderías, seguros marítimos, arbitraje comercial internacional y derecho anglosajón aplicable por elección de las partes. Reconocer términos transparentes, identificar falsos cognados (como \emph{jurisprudence} o \emph{statute}) y comprender cómo se organiza un sistema jurídico distinto del propio evita errores de interpretación que, en la práctica profesional, pueden significar un contrato mal entendido, una cláusula mal aplicada o un dictamen erróneo. Además, comprender por qué el derecho anglosajón privilegia la estabilidad y la tradición explica por qué gran parte del comercio internacional elige someterse a esa legislación.
\end{note}

\section{Utilidad general de la herramienta de lectocomprensión}

\begin{note}
La estrategia trabajada en clase --- identificar la estructura general de un texto antes que el detalle, apoyarse en palabras reconocibles y desconfiar de aquellas que parecen conocidas pero no lo son --- constituye una herramienta transferible a cualquier situación en la que deba comprenderse un texto técnico en un idioma que no se domina por completo. Es una habilidad de lectura crítica que sirve tanto para leer un artículo académico o un manual técnico, como para comparar la propia realidad institucional con la de otro país y comprenderla mejor por contraste.
\end{note}

\section{Relación entre los bloques temáticos}

\begin{note}
Puede pensarse la organización de la clase como un edificio. Antes de entrar a mirar las habitaciones (los detalles del texto), conviene mirar el plano general. Recién entonces corresponde preguntar qué es una Constitución ``en general'', para descubrir que la británica está construida con materiales distintos a los del sistema continental: sin un plano único, sino con piezas sueltas que son las costumbres, las leyes del Parlamento y los fallos judiciales. Varias de esas palabras en inglés (\emph{jurisprudence}, \emph{statute}, \emph{convention}) parecen ``piezas conocidas'' pero son en realidad piezas distintas: son los falsos cognados que pueden hacer que la lectura se derrumbe si no se identifican a tiempo. Con esos materiales identificados, puede comprenderse quién sostiene el edificio (la soberanía parlamentaria y la Corona) y qué reformas fue teniendo con el paso del tiempo.
\end{note}

% ===========================================================
\chapter{Metodología de lectura de textos jurídicos en inglés}
% ===========================================================

\section{Presentación del texto y su fuente}

El texto de trabajo, identificado en clase como ``texto 3'', resulta más extenso que los materiales trabajados previamente y, a diferencia de otros textos ya analizados, no presenta subtítulos ni una división visual marcada más allá de los párrafos. Por ello, la primera lectura debe apoyarse en otros indicios: el título, la fuente del texto y los términos reconocibles.

\section{Identificación de términos transparentes}

Como estrategia central de lectocomprensión, se trabajó sobre la identificación de los términos transparentes del texto, es decir, aquellas palabras cuya semejanza con el español permite anticipar su significado sin necesidad de traducción literal.

\begin{example}{Términos transparentes identificados en el texto}
Palabras como \emph{Constitution}, \emph{organize}, \emph{distribute} y \emph{regulate} remiten directamente a los conceptos en español de Constitución, organizar, distribuir y regular. Ya desde el título (\emph{Constitution}) y las primeras líneas del texto, el reconocimiento de este vocabulario permite anticipar el tema general sin necesidad de traducir palabra por palabra.
\end{example}

\begin{note}
No todos los términos relevantes son transparentes por semejanza ortográfica: algunos se infieren por asociación con otras palabras conocidas. Así ocurre con el término \emph{unwritten}: aunque \emph{written} no se parece ortográficamente a ``escrito'', el reconocimiento de la raíz \emph{writing} permite inferir su significado (lo no escrito, lo no codificado).
\end{note}

% ===========================================================
\chapter{Contexto geográfico y político: Reino Unido e Inglaterra}
% ===========================================================

Antes de avanzar sobre el contenido jurídico del texto, se precisó una cuestión de cultura general necesaria para evitar un error frecuente.

\begin{important}
Un error común consiste en identificar ``Reino Unido'' con ``Inglaterra''. El nombre oficial del país es \emph{United Kingdom of Great Britain and Northern Ireland}. \emph{Great Britain} abarca Inglaterra, Gales y Escocia. Inglaterra es, por lo tanto, solo una parte del Reino Unido, aunque el error es frecuente porque allí se encuentra la sede del gobierno.
\end{important}

\begin{note}
Irlanda del Norte forma parte del Reino Unido, mientras que Irlanda (del Sur) queda por fuera de esta unión política. Se trata de una precisión de cultura general y geográfica que no depende del texto, pero que resulta indispensable para comprenderlo correctamente.
\end{note}

% ===========================================================
\chapter{La Constitución: definición general y particularidad del Reino Unido}
% ===========================================================

\section{Qué es una Constitución, en términos generales}

El primer párrafo del texto comienza con una definición general aplicable a cualquier Constitución, para luego contrastarla con el caso particular del Reino Unido.

\begin{formula}{Definición general de Constitución}
\emph{``Constitutions organize, distribute and regulate state power.''}

Traducción de trabajo consensuada en la clase: las constituciones organizan, distribuyen y regulan el poder estatal.
\end{formula}

\begin{definition}{Tres funciones de toda Constitución}
\begin{itemize}
    \item \textbf{Organizar} el poder estatal: establece la estructura del Estado y sus principales instituciones (\emph{major state institutions}), es decir, aquellas establecidas en la Constitución --- como la división tripartita de poderes en Estados Unidos o la existencia de la Corona en el caso británico.
    \item \textbf{Distribuir} el poder estatal: en el caso de Estados Unidos, mediante la división en tres poderes; en el caso del Reino Unido, con la particularidad de la existencia de la Corona.
    \item \textbf{Regular} las relaciones entre esas instituciones (\emph{with each other}) y entre ellas y la ciudadanía (\emph{with the state citizens}): los principios que rigen o controlan (\emph{governing}) esos vínculos.
\end{itemize}
\end{definition}

\section{``Unwritten'' y ``uncodified'': una distinción clave}

El texto original utiliza el término \emph{unwritten} entre comillas, lo cual no es casual, sino una señal para el lector: el término resulta impreciso, porque sí existen normas escritas dentro del ordenamiento británico (las leyes del Parlamento, entre otras). Lo que en rigor falta no es la escritura, sino la codificación en un único cuerpo normativo.

\begin{important}
\textbf{No escrita $\neq$ no codificada.} El Reino Unido sí cuenta con normas escritas; lo que no tiene es un único instrumento donde estén codificadas, a diferencia de sistemas como el argentino o el estadounidense, en los que existe un documento único al que puede acudirse para encontrar el contenido constitucional.
\end{important}

\section{Ventajas y desventajas de la dispersión normativa}

\begin{note}
\textbf{Ventaja:} mayor facilidad para modificar el sistema, precisamente por no estar concentrado en un único instrumento formal que exija un procedimiento de reforma agravado.

\textbf{Desventaja:} mayor dificultad de acceso y sistematización para quienes deben invocar o fundamentar un principio constitucional, a diferencia de lo que ocurre en un sistema como el argentino, donde puede acudirse directamente al texto de la Constitución.
\end{note}

% ===========================================================
\chapter{Origen histórico y estabilidad del sistema constitucional británico}
% ===========================================================

\section{Evolución frente a revolución}

\begin{example}{Contraste señalado en el texto}
\emph{``[\ldots] By contrast, the British Constitution has evolved over a long period of time [\ldots] reflecting the relative stability of the British polity.''}

El conector \emph{by contrast} marca la diferencia entre los países que debieron ``empezar de cero'' (\emph{start from scratch}) tras un cambio de régimen --- como la Argentina tras su independencia de España, que debió constituir un régimen nuevo prácticamente desde cero --- y el caso británico, de evolución gradual y continua, sin una ruptura institucional equivalente.
\end{example}

Esta ausencia de ruptura institucional explica, según se señaló en clase, la estabilidad relativa del sistema político británico (\emph{the relative stability of the British polity}).

\section{Estabilidad británica y comercio internacional}

\begin{example}{Proyección de la estabilidad británica en el comercio internacional}
El derecho mercantil vinculado a la compraventa internacional de mercaderías se rige, en gran medida, por el \emph{Sale of Goods Act}. La regulación internacional de los contratos elige con frecuencia el derecho británico porque este se ha mantenido sustancialmente estable desde el siglo XV. En materia de seguros marítimos, la tradición de \emph{Lloyd's} --- la primera compañía de seguros marítimos y de cargamento --- constituye un antecedente que el resto del comercio internacional, incluido Estados Unidos, terminó adoptando.
\end{example}

\begin{note}
El Reino Unido es una cultura organizada en función de la tradición y del respeto a esta, a diferencia de sistemas como el argentino, donde los cambios de gobierno y de orientación política suelen implicar modificaciones significativas en la estructura estatal.
\end{note}

% ===========================================================
\chapter{Falsos cognados jurídicos fundamentales}
% ===========================================================

Uno de los ejes centrales de la clase fue la identificación de falsos cognados (\emph{false friends}) presentes en el texto, es decir, palabras que por su semejanza ortográfica con el español inducen a una traducción incorrecta.

\section{\emph{Jurisprudence} no es jurisprudencia}

\begin{important}
El término inglés \emph{jurisprudence} no significa ``jurisprudencia'' en el sentido de fallos judiciales, sino que hace referencia a la filosofía del derecho o a la doctrina, en un sentido más próximo a la teoría del derecho. Para referirse a la jurisprudencia (en el sentido de fallos judiciales) se emplean en cambio los términos \emph{legal precedent}, \emph{case law} o \emph{judicial decisions}.
\end{important}

\section{\emph{Statute} no es estatuto}

\begin{important}
En el derecho inglés, \emph{statute} no equivale a ``estatuto'' en el sentido de norma interna de una institución u organización (como un estatuto universitario o el estatuto social de una sociedad comercial). Un \emph{statute} es una ley formal y escrita, sancionada por el Parlamento (el órgano legislativo), equivalente a lo que en español se llama simplemente ``ley''.
\end{important}

Esta distinción se vincula con la pirámide de fuentes del derecho: en el sistema continental, la jurisprudencia ocupa un lugar secundario frente a la Constitución y los tratados, mientras que en el sistema anglosajón el \emph{common law} se ubica en un escalón mucho más alto, dada su relevancia como fuente autónoma de derecho.

\section{\emph{Convention} no es convención internacional}

\begin{important}
Una \emph{convention}, en el sistema constitucional británico, no es un tratado ni una convención internacional. Se trata de prácticas no escritas (\emph{unwritten practices}), no codificadas, que se fueron desarrollando con el tiempo y que regulan \emph{the business of governing}: el ejercicio de gobernar y administrar el Estado. Es, en definitiva, costumbre institucional consolidada por la tradición.
\end{important}

\section{Cuadro comparativo de falsos cognados}

\begin{table}[H]
\centering
\caption{Falsos cognados jurídicos identificados en el texto}
\begin{tabularx}{\textwidth}{
    >{\raggedright\arraybackslash}p{0.22\textwidth}
    >{\raggedright\arraybackslash}p{0.33\textwidth}
    >{\raggedright\arraybackslash}X
}
\toprule
\textbf{Término en inglés} & \textbf{Falso sentido / traducción intuitiva} & \textbf{Significado correcto en derecho} \\
\midrule
\emph{Jurisprudence} & Jurisprudencia & Filosofía del derecho / doctrina \\
\emph{Statute} & Estatuto (social, universitario) & Ley formal escrita, sancionada por el Parlamento \\
\emph{Convention} & Convención internacional / tratado & Costumbre o práctica institucional no escrita \\
\bottomrule
\end{tabularx}
\end{table}

% ===========================================================
\chapter{Fuentes de la Constitución británica}
% ===========================================================

\section{La ``pirámide de Kelsen forzada'' como recurso didáctico}

\begin{note}
Para visualizar el cambio en la jerarquía de fuentes entre el sistema continental y el anglosajón, en clase se propuso adaptar gráficamente la pirámide de Kelsen, aclarando expresamente su carácter didáctico y no doctrinario: se trata de un recurso propio de la clase, no de una construcción reconocida por la doctrina constitucionalista, cuya única finalidad es visualizar cómo cambia el orden de los elementos según el sistema jurídico de referencia.
\end{note}

\begin{table}[H]
\centering
\caption{Comparación de la jerarquía de fuentes: sistema continental vs. sistema anglosajón}
\begin{tabularx}{\textwidth}{
    >{\raggedright\arraybackslash}p{0.24\textwidth}
    >{\raggedright\arraybackslash}p{0.36\textwidth}
    >{\raggedright\arraybackslash}X
}
\toprule
\textbf{Nivel de la pirámide} & \textbf{Sistema continental} & \textbf{Sistema anglosajón (Reino Unido)} \\
\midrule
Nivel superior & Constitución y tratados con jerarquía constitucional & \emph{Statutes}: leyes del Parlamento \\
Nivel intermedio & Leyes ordinarias & \emph{Conventions}: costumbre institucional no escrita \\
Nivel inferior / relevante & Jurisprudencia y costumbre, con escasa relevancia como fuente formal & \emph{Common law} / \emph{case law}: decisiones judiciales, con altísima relevancia como fuente \\
Nivel más bajo & --- & Doctrina / \emph{constitutional authorities}, de carácter consultivo \\
\bottomrule
\end{tabularx}
\end{table}

\section{Enumeración de las fuentes según el texto}

\begin{formula}{Fuentes enumeradas en el texto de trabajo}
\emph{``The British Constitution [is] derived from a number of sources [\ldots] statutes, conventions, judicial decisions [\ldots]''}

A lo que se suman, según el desarrollo posterior del texto, el derecho de la Unión Europea (con las salvedades del Brexit), el derecho internacional y las \emph{constitutional authorities}.
\end{formula}

\subsection*{a) \emph{Statutes} --- leyes del Parlamento}

Los \emph{statutes} son las normas que emanan del Parlamento: \emph{``laws passed by Parliament [\ldots] and the highest form of law''}. Constituyen la fuente jerárquicamente más alta dentro de esta pirámide adaptada.

\subsection*{b) \emph{Conventions} --- costumbre institucional}

Las \emph{conventions} constituyen prácticas no escritas que regulan el ejercicio del gobierno, consolidadas por tradición: determinadas potestades quedan reservadas por costumbre a ciertas instituciones (por ejemplo, al Parlamento), sin que exista una norma escrita que así lo establezca.

\subsection*{c) \emph{Common law} --- decisiones judiciales}

\begin{example}{Explicación del \emph{common law}}
El \emph{common law} es \emph{``law developed by the courts and judges through cases''}: la ley común o jurisprudencia constituida por los fallos judiciales. Se trata de la aplicación por analogía de la solución adoptada en un caso a otras causas con identidad de circunstancias.
\end{example}

\subsection*{d) Derecho de la Unión Europea y el impacto del Brexit}

El texto de trabajo resulta, en este punto, desactualizado, dado que hace referencia a la incorporación de normas europeas al ordenamiento británico, un proceso hoy superado tras la salida del Reino Unido de la Unión Europea. Antes del Brexit, el Reino Unido debió incorporar normas provenientes de la Comunidad Europea, estableciéndose un plazo para resolver qué normas conservar dentro de su ordenamiento interno y cuáles no.

\subsection*{e) Derecho internacional}

El texto agrega que el Reino Unido está también sujeto al derecho internacional (\emph{UK is also subject to international law}), lo cual se vincula con la imposibilidad de aislamiento normativo de cualquier Estado moderno respecto de sus compromisos internacionales.

\subsection*{f) \emph{Constitutional authorities} --- autoridades constitucionales y doctrina}

\begin{note}
Ante la ausencia de un texto único al que recurrir, cumplen una función consultiva ciertas autoridades y especialistas --- jueces, abogados y doctrinarios consultados como \emph{constitutional authorities}. Se trata de una función asimilable a una recomendación, razón por la cual se ubica en el nivel más bajo de la pirámide comparada, a diferencia de la jurisprudencia y la costumbre, que sí tienen fuerza normativa directa en este sistema.
\end{note}

\begin{note}
Como ejemplo institucional adicional, la Cámara de los Lores (\emph{House of Lords}) constituye la cámara alta del Parlamento británico. Hasta hace relativamente poco tiempo, la mayoría de sus bancas se transmitían de manera hereditaria, ligadas a la nobleza; actualmente sus integrantes son elegidos, en una reforma que ilustra el carácter fuertemente tradicional de esta estructura gubernamental.
\end{note}

% ===========================================================
\chapter{Soberanía parlamentaria y principios centrales}
% ===========================================================

\section{La Corona y el Parlamento: origen de la potestad legislativa}

El texto sintetiza la organización del sistema británico en pocas palabras.

\begin{formula}{Síntesis textual del principio organizador}
\emph{``What the Queen in Parliament enacts is law.''}

Traducción de trabajo: lo que la Corona en el Parlamento promulga es ley. El texto utiliza \emph{``the Queen''} por tratarse de una versión anterior al cambio de titularidad de la Corona británica.
\end{formula}

El Parlamento posee potestades y facultades porque se las otorga la Corona: se trata de una facultad delegada, y el Parlamento promulga las leyes \emph{using the power of the Crown}. Esto se diferencia del sistema argentino, donde la potestad legislativa corresponde al Congreso porque se la otorga la Constitución; en el Reino Unido, en cambio, es la Corona quien históricamente la otorga al Parlamento. De allí deriva el concepto de \emph{law-making power} (poder de creación de leyes): los legisladores son, literalmente, ``hacedores de leyes'' (\emph{law-makers}).

\begin{definition}{Soberanía parlamentaria (\emph{parliamentary sovereignty})}
\emph{``Parliamentary sovereignty is regarded as the defining principle of the British Constitution [\ldots] this is the ultimate law-making power, vested in a democratically elected Parliament, to create or abolish any law.''}

Es el principio definitorio de la Constitución británica según el cual el Parlamento, democráticamente electo, posee el poder legislativo último (\emph{ultimate law-making power}) para crear o abolir cualquier ley. El par \emph{create or abolish} expresa la amplitud de esta potestad; el carácter ``último'' no remite a un orden jerárquico dentro de un ordenamiento, sino al carácter definitivo de esa potestad, que recae en el centro del sistema unitario: el Parlamento.
\end{definition}

\section{Principios centrales de la Constitución británica}

\begin{definition}{Principios centrales identificados en el texto}
\begin{itemize}
    \item \textbf{The rule of law}: el estado de derecho o ``imperio de la ley''. No debe traducirse literalmente como ``regla de la ley''; su sentido jurídico correcto remite al concepto de estado de derecho.
    \item \textbf{The separation of executive, legislative and judicial branches}: la separación entre los poderes ejecutivo, legislativo y judicial.
    \item La existencia de un \textbf{estado unitario} (\emph{a unitary state}), en el cual el poder último o definitivo (\emph{ultimate power}) recae en el centro: el Parlamento.
\end{itemize}
\end{definition}

\section{El carácter ``mítico'' de estos principios y las tensiones actuales}

El texto introduce, mediante el conector \emph{however} (sin embargo), un matiz crítico respecto de los principios anteriormente descriptos.

\begin{important}
\emph{``However, some of these principles are [\ldots] mythical. The British Constitution may be better [described] as the fusion of executive and legislative [\ldots] and Parliamentary sovereignty may now be [\ldots] the impact of Europe [\ldots]''}

Pese a proclamarse la separación de poderes, en la práctica existe una fusión entre el ejecutivo y el legislativo, sostenida por la tradición histórica más que por una estructura de Estado moderno estrictamente diferenciada. El término \emph{mythical} no remite a algo inexistente, sino a una construcción sostenida por la tradición histórica que hoy no encaja del todo con la estructura de un Estado moderno.
\end{important}

\section{Impacto del derecho europeo sobre la soberanía parlamentaria}

El cuestionamiento a la soberanía parlamentaria como principio absoluto surgió, entre otros factores, a partir de la incorporación de normativa proveniente de la Comunidad Europea, previa al Brexit: el Parlamento debió incorporar legislación, tratados y jurisprudencia provenientes de instituciones europeas ajenas a su propia estructura soberana, generando una tensión que se resolvió parcialmente tras el Brexit.

\section{Vigencia y legitimidad actual de la Corona}

\begin{note}
Se encuentra en discusión, en los últimos años, el costo que representa para la ciudadanía mantener a los integrantes de la Corona y de la nobleza, así como el rol que estos cumplen en la actualidad. Como ejemplo histórico adicional, hacia el año 2015 se modificaron las reglas de sucesión al trono, que hasta entonces otorgaban preferencia al hijo varón sobre las hijas mujeres nacidas con anterioridad.
\end{note}

% ===========================================================
\chapter{Problemas y reformas de una Constitución no codificada}
% ===========================================================

\section{Dos problemas derivados de la falta de codificación}

\begin{formula}{Planteo textual}
\emph{``A[n] uncodified constitution creates two problems.''}
\end{formula}

\begin{important}
\textbf{Primer problema: falta de certeza.} Al encontrarse la Constitución ``desparramada'' en múltiples fuentes, resulta difícil identificar dónde reside el contenido constitucional aplicable a un caso, a diferencia de lo que ocurre en sistemas como el argentino o el estadounidense, donde puede acudirse directamente al texto codificado.
\end{important}

\begin{important}
\textbf{Segundo problema: mayor facilidad de reforma.} Los sistemas con constituciones escritas cuentan con \emph{documents with higher law}: un documento con estatus de ley superior, que solo puede modificarse mediante un procedimiento agravado (\emph{amendable elaborate procedure}). Al no existir en el Reino Unido un instrumento codificado equivalente, el procedimiento de modificación es, en gran medida, la propia costumbre.
\end{important}

\begin{example}{El Brexit como modificación constitucional}
El Brexit constituye un ejemplo concreto de esta flexibilidad: al salir de la Comunidad Europea, el Reino Unido debió analizar qué normas europeas quería mantener y cuáles dejar fuera de su ordenamiento, sin que dicha modificación quedara plasmada en un instrumento constitucional único.
\end{example}

\section{Reformas históricas desde 1997}

\begin{formula}{Reformas enumeradas en el texto}
\emph{``The flexibility of the UK constitution is evident [since 1997], including the abolition of the majority of hereditary [seats] in the House of Lords.''}
\end{formula}

Se trata de reformas que, al no requerir un procedimiento agravado único, se fueron incorporando de manera sucesiva y dispersa:

\begin{itemize}
    \item La abolición de la mayoría de las bancas hereditarias en la Cámara de los Lores (\emph{House of Lords}).
    \item La incorporación del \emph{Human Rights Act} (ley de derechos humanos).
    \item Los \emph{devolution settlements} otorgados a Escocia, Gales e Irlanda del Norte.
\end{itemize}

\section{\emph{Devolution settlements}: delegación revocable de potestades}

\begin{definition}{\emph{Devolution settlements}}
Son pactos mediante los cuales el gobierno central del Reino Unido delega potestades propias a los gobiernos de Escocia, Gales e Irlanda del Norte. A diferencia de lo que ocurre en un sistema federal, esta delegación es revocable: \emph{``this devolution settlement could ever be repealed''}, es decir, podría en teoría ser derogada por el propio Reino Unido, que podría decidir no delegar más esas facultades.
\end{definition}

\begin{note}
Este mecanismo se vincula con la actualidad política del Reino Unido, en la que existen tensiones recurrentes entre el gobierno central y Escocia, Gales e Irlanda del Norte respecto de la extensión y reversibilidad de las potestades delegadas.
\end{note}

\section{Precisión terminológica final: \emph{British}, \emph{English} y \emph{UK}}

\begin{note}
El gentilicio para el Reino Unido en su conjunto es \emph{British}; no existe otro gentilicio distinto. Según el nombre oficial, Gran Bretaña comprende Inglaterra, Gales y Escocia, mientras que Irlanda del Norte es un territorio distinto; sin embargo, al utilizar el gentilicio \emph{British} se entiende que este abarca a todo el Reino Unido. Se trata de un fenómeno comparable al uso de ``americano'' para referirse exclusivamente a los ciudadanos de Estados Unidos, pese a que ``América'' designa a todo el continente: un caso de apropiación lingüística de un término colectivo por parte del país dominante, siendo el gentilicio técnicamente correcto, en ese caso, ``estadounidense''.
\end{note}

% ===========================================================
\chapter{Cuadro Cornell de repaso}
% ===========================================================

A continuación se presenta el cuadro de repaso organizado según el método Cornell: preguntas y conceptos clave a la izquierda, desarrollo y explicación a la derecha, con una síntesis integradora al final.

\begin{longtable}{
    >{\raggedright\arraybackslash}p{0.30\textwidth}
    >{\raggedright\arraybackslash}p{0.60\textwidth}
}
\caption{Cuadro Cornell --- ``What is the UK Constitution?''} \\
\toprule
\textbf{Preguntas / palabras clave} & \textbf{Notas / respuestas} \\
\midrule
\endfirsthead

\multicolumn{2}{c}{\textit{(continúa de la página anterior)}} \\
\toprule
\textbf{Preguntas / palabras clave} & \textbf{Notas / respuestas} \\
\midrule
\endhead

\midrule
\multicolumn{2}{r}{\textit{(continúa en la página siguiente)}} \\
\endfoot

\bottomrule
\endlastfoot

\textbf{¿Qué tres funciones cumple toda Constitución, según el texto?} &
Organizar, distribuir y regular el poder estatal (\emph{organize, distribute and regulate state power}): establece las principales instituciones del Estado y los principios que rigen las relaciones entre ellas y con la ciudadanía. \\

\textbf{¿Cuál es la diferencia entre ``unwritten'' y ``uncodified''?} &
\emph{Unwritten} (no escrita) es impreciso, porque sí existen normas escritas en el Reino Unido (leyes del Parlamento). Lo correcto es \emph{uncodified} (no codificada): las normas existen, pero no están reunidas en un único instrumento. \\

\textbf{¿Por qué el Reino Unido no tuvo que ``empezar de cero'' como otros países?} &
Porque no atravesó una revolución ni un cambio de régimen que exigiera fundar nuevas instituciones; su sistema evolucionó gradualmente a lo largo de un período extenso, lo que explica su estabilidad relativa (\emph{relative stability of the British polity}). \\

\textbf{¿Por qué se elige frecuentemente el derecho inglés en los contratos internacionales?} &
Por su estabilidad histórica: el derecho comercial y marítimo inglés (por ejemplo, el \emph{Sale of Goods Act} y la tradición de \emph{Lloyd's} en seguros marítimos) se mantiene sustancialmente estable desde el siglo XV, lo que brinda previsibilidad a las partes contratantes. \\

\textbf{¿Qué significa \emph{jurisprudence} en inglés jurídico y por qué es un falso cognado?} &
Significa filosofía del derecho o doctrina, no ``jurisprudencia'' en el sentido de fallos judiciales. Para referirse a la jurisprudencia en español se usan los términos \emph{case law}, \emph{legal precedent} o \emph{judicial decisions}. \\

\textbf{¿Qué es un \emph{statute} en el derecho inglés?} &
Una ley formal y escrita sancionada por el Parlamento (el órgano legislativo), equivalente a lo que en español llamamos simplemente ``ley''. No debe confundirse con ``estatuto'' en el sentido de norma interna de una institución o sociedad. \\

\textbf{¿Qué son las \emph{conventions} en el sistema constitucional británico?} &
Prácticas institucionales no escritas que regulan el ejercicio del gobierno (\emph{the business of governing}), consolidadas por la costumbre y la tradición a lo largo del tiempo. No deben confundirse con ``convención'' en el sentido de tratado internacional. \\

\textbf{¿Cuáles son las principales fuentes de la Constitución británica?} &
Los \emph{statutes} (leyes del Parlamento), las \emph{conventions} (costumbre institucional), el \emph{common law} (decisiones judiciales), el derecho de la Unión Europea (residual, tras el Brexit), el derecho internacional y las \emph{constitutional authorities} (doctrina, de carácter consultivo). \\

\textbf{¿Cómo cambia la jerarquía de fuentes entre el sistema continental y el anglosajón?} &
En el sistema continental, la jurisprudencia y la costumbre ocupan un lugar bajo en la pirámide de fuentes, por debajo de la Constitución y los tratados. En el sistema anglosajón, el \emph{common law} y la costumbre (\emph{conventions}) ocupan un lugar mucho más relevante, casi en la cúspide. \\

\textbf{¿De dónde deriva la potestad legislativa del Parlamento británico?} &
De una facultad delegada por la Corona: \emph{``what the Queen in Parliament enacts is law''}. A diferencia del sistema argentino, donde la potestad legislativa emana directamente de la Constitución, en el Reino Unido es la Corona quien históricamente la otorga al Parlamento. \\

\textbf{¿Qué es la soberanía parlamentaria (\emph{parliamentary sovereignty})?} &
El principio definitorio de la Constitución británica según el cual el Parlamento, democráticamente electo, posee el poder legislativo último (\emph{ultimate law-making power}) para crear o abolir cualquier ley (\emph{to create or abolish any law}). \\

\textbf{¿Qué otros principios centrales identifica el texto, además de la soberanía parlamentaria?} &
El \emph{rule of law} (estado de derecho), la separación entre los poderes ejecutivo, legislativo y judicial, y la existencia de un Estado unitario en el que el poder último recae en el Parlamento como centro. \\

\textbf{¿Por qué se dice que algunos principios de la Constitución británica son ``míticos''?} &
Porque, pese a proclamarse la separación de poderes, en la práctica existe una fusión entre el ejecutivo y el legislativo, sostenida por la tradición histórica más que por una estructura de Estado moderno estrictamente diferenciada. \\

\textbf{¿Qué impacto tuvo el derecho de la Unión Europea sobre la soberanía parlamentaria?} &
Cuestionó su carácter absoluto, dado que el Parlamento debió incorporar legislación, tratados y jurisprudencia provenientes de instituciones europeas ajenas a su propia estructura soberana, generando una tensión que se resolvió parcialmente tras el Brexit. \\

\textbf{¿Qué dos problemas genera una Constitución no codificada?} &
En primer lugar, una falta de certeza sobre dónde encontrar el contenido constitucional aplicable a un caso. En segundo lugar, una mayor facilidad para reformarla, al no exigir un procedimiento agravado único como en los sistemas con constituciones escritas y de jerarquía superior. \\

\textbf{¿Qué son los \emph{devolution settlements}?} &
Acuerdos mediante los cuales el gobierno central del Reino Unido delega potestades propias a los gobiernos de Escocia, Gales e Irlanda del Norte. A diferencia de un sistema federal, esta delegación es revocable: podría, en teoría, ser derogada (\emph{repealed}) por el propio Reino Unido. \\

\textbf{¿Por qué se usa ``British'' como gentilicio único para todo el Reino Unido?} &
Porque no existe otro gentilicio distinto para referirse a los ciudadanos del Reino Unido en su conjunto, pese a que, técnicamente, Gran Bretaña (Inglaterra, Gales y Escocia) e Irlanda del Norte son territorios distintos. Es un fenómeno comparable al uso de ``americano'' para referirse solo a Estados Unidos. \\

\multicolumn{2}{p{0.90\textwidth}}{
\vspace{0.3em}
\textbf{Resumen:} La clase analizó el texto \emph{``What is the UK constitution?''} como caso de estudio para desarrollar estrategias de lectocomprensión en inglés jurídico --- identificación de macroestructura y de términos transparentes --- y, al mismo tiempo, para comprender en profundidad un sistema constitucional radicalmente distinto del propio. El Reino Unido no carece de Constitución, como suele creerse erróneamente, sino que la tiene dispersa en múltiples fuentes no codificadas: las leyes del Parlamento (\emph{statutes}), la costumbre institucional (\emph{conventions}), los precedentes judiciales (\emph{common law}), el derecho internacional y, residualmente, el derecho europeo, todo ello complementado por la función consultiva de la doctrina (\emph{constitutional authorities}). Esta dispersión, lejos de ser un simple defecto técnico, expresa un modo de organización estatal fundado en la evolución histórica y la tradición, antes que en una ruptura fundacional como la vivida por Argentina o Estados Unidos, y explica tanto la estabilidad que hace atractivo al derecho inglés en el comercio internacional como los dos grandes problemas que genera: la incertidumbre sobre dónde reside el contenido constitucional y la facilidad relativa con que puede modificarse, como se evidenció con las reformas posteriores a 1997 y con el propio Brexit. A lo largo del recorrido, se puso especial énfasis en tres falsos cognados --- \emph{jurisprudence}, \emph{statute} y \emph{convention} --- que, de no identificarse correctamente, llevarían a una comprensión distorsionada tanto del texto como del sistema jurídico que describe, subrayando que la lectocomprensión de textos jurídicos exige no solo reconocer el vocabulario transparente, sino también desconfiar de aquel que, por parecido, puede inducir a error.
\vspace{0.3em}
} \\

\end{longtable}

\end{document}
